\documentclass[11pt]{article}
\usepackage{amsmath,amssymb}
\title{ROOT V142: Helix, Pure Gauge and Photon Polarization Audit}
\author{Reproducible Mathematical Research Notes}
\begin{document}
\maketitle
\paragraph{Units.} $c=1$ only; $\hbar$ and $G$ retain dimensions.
\section{Helical kinematics}
\begin{equation}\mathbf r(t)=(a\cos\omega t,a\sin\omega t,ut),\quad a^2\omega^2+u^2=1.\end{equation}
\begin{equation}\kappa=a\omega^2,\quad\tau=u\omega,\quad\omega^2=\kappa^2+\tau^2.\end{equation}
\section{Flat full rotational connection}
\begin{equation}\Omega=R^{-1}dR,\quad d\Omega+\Omega\wedge\Omega=0,\qquad R:U\to SO(3).\end{equation}
Projected curvature can be nonzero, e.g. with $\theta=\alpha x,\phi=\beta y$:
\begin{equation}f_{xy}=\tfrac12\mathbf n\cdot(\partial_x\mathbf n\times\partial_y\mathbf n)=\tfrac12\alpha\beta\sin(\alpha x).\end{equation}
\section{Correct circularly polarized Maxwell wave}
\begin{align}
\varphi&=k(t-z),\quad k>0,\\
\mathbf E&=E_0(\cos\varphi,\sin\varphi,0),\\
\mathbf B&=E_0(-\sin\varphi,\cos\varphi,0),\\
\partial_t\mathbf B+\nabla\times\mathbf E&=0,\quad\partial_t\mathbf E-\nabla\times\mathbf B=0,\\
\mathbf E\times\mathbf B&=E_0^2\hat{\mathbf z}.
\end{align}
The opposite sign for $\mathbf B$ is not a vacuum solution with this phase.
\section{Gauge constraints and propagating modes}
\begin{equation}N_{\rm Maxwell}=\frac{8-2\times2}{2}=2,\qquad N_{\rm Proca}=\frac{8-2}{2}=3.\end{equation}
\section{Constant background composite field no-go}
\begin{equation}
q=q_0+\varepsilon\xi,\quad p=p_0+\varepsilon\eta,\quad
F=dq\wedge dp=O(\varepsilon^2),\quad S_F=O(\varepsilon^4).
\end{equation}
\section{Dimensional consistency}
\begin{equation}
[\hbar]=ML,\quad[G]=L/M,\quad [F]=L^{-2},\quad [R]=L^{-2}.
\end{equation}
\begin{equation}
S=\int d^4x\sqrt{-g}\left(\frac{R}{16\pi G}-\frac{\hbar}{4g^2}F_{\mu\nu}F^{\mu\nu}\right),
\end{equation}
which reproduces known theories by assumption, not a demonstrated unification.
\end{document}
